% Standalone problem document, generated from the adjacent README.md.
% Compile with XeLaTeX. Mathematical content is maintained in Markdown.
\documentclass[11pt,a4paper]{article}
\usepackage[fontsize=10.5pt]{scrextend}
\usepackage[margin=22mm,headheight=15pt,headsep=9mm,footskip=11mm]{geometry}
\usepackage{fontspec}
\setmainfont{texgyrepagella}[Extension=.otf,UprightFont=*-regular,BoldFont=*-bold,ItalicFont=*-italic,BoldItalicFont=*-bolditalic]
\setsansfont{texgyreheros}[Extension=.otf,UprightFont=*-regular,BoldFont=*-bold,ItalicFont=*-italic,BoldItalicFont=*-bolditalic]
\setmonofont{texgyrecursor}[Extension=.otf,UprightFont=*-regular,BoldFont=*-bold,ItalicFont=*-italic,BoldItalicFont=*-bolditalic,Scale=MatchLowercase]
\usepackage{amsmath,amssymb}
\usepackage{unicode-math}
\setmathfont{texgyrepagella-math.otf}
\usepackage{xcolor,microtype,enumitem,needspace,fancyhdr}
\usepackage{bookmark}
\definecolor{ink}{HTML}{17344B}
\definecolor{accent}{HTML}{197D80}
\definecolor{muted}{HTML}{596773}
\hypersetup{colorlinks=true,linkcolor=accent,urlcolor=accent,pdftitle={RA-09 - Concave-function
transfer of Frobenius low-rank
error},pdfauthor={Open Problems in Numerical Linear Algebra}}
\urlstyle{same}
\setlength{\parindent}{0pt}
\setlength{\parskip}{5pt plus 1pt minus 1pt}
\linespread{1.04}
\setlength{\emergencystretch}{3em}
\widowpenalty=10000
\clubpenalty=10000
\displaywidowpenalty=10000
\allowdisplaybreaks[1]
\setlist{nosep,leftmargin=1.5em}
\providecommand{\tightlist}{\setlength{\itemsep}{0pt}\setlength{\parskip}{0pt}}
\providecommand{\pandocbounded}[1]{#1}
\pagestyle{fancy}
\fancyhf{}
\fancyhead[L]{\sffamily\footnotesize\color{muted}OPEN PROBLEMS IN NUMERICAL LINEAR ALGEBRA}
\fancyhead[R]{\sffamily\bfseries\footnotesize\color{accent}RA-09}
\fancyfoot[L]{\parbox[b]{0.9\textwidth}{\sffamily\fontsize{8}{10}\selectfont\color{muted}Editorial ratings; a bounded search cannot certify that a problem remains unsolved.\\Literature check: 2026-09-11}}
\fancyfoot[R]{\sffamily\footnotesize\color{muted}\thepage}
\renewcommand{\headrulewidth}{0.3pt}
\renewcommand{\headrule}{\hbox to\headwidth{\color{accent}\leaders\hrule height \headrulewidth\hfill}}
\makeatletter
\renewcommand{\section}{\Needspace{5\baselineskip}\@startsection{section}{1}{0pt}{12pt plus 2pt minus 2pt}{4pt}{\sffamily\bfseries\large\color{ink}}}
\renewcommand{\subsection}{\Needspace{4\baselineskip}\@startsection{subsection}{2}{0pt}{9pt plus 1pt minus 1pt}{3pt}{\sffamily\bfseries\normalsize\color{ink}}}
\makeatother
\setcounter{secnumdepth}{0}
\begin{document}
{\sffamily\small\color{muted}Randomized and low-rank approximation\par}
\vspace{3pt}
{\raggedright\hyphenpenalty=10000\exhyphenpenalty=10000\sffamily\bfseries\fontsize{21}{24}\selectfont\color{ink}Concave-function
transfer of Frobenius low-rank error\par}
\vspace{7pt}
{\sffamily\small\textbf{Difficulty:} challenging\quad\textbf{Importance:} interesting
to the community\par}
{\sffamily\small\bfseries\color{accent}Status: Solved\par}
\vspace{4pt}
{\color{accent}\hrule height 0.8pt}
\vspace{6pt}
\textbf{Rating rationale:} Challenging because squared Frobenius
transfer needs control of eigenvectors as well as eigenvalues; community
impact is reusable low-rank approximation of matrix functions.\\
\textbf{Topic:} Low-rank approximation of matrix functions.

\subsection{Resolution --- 2026-09-11}\label{resolution-2026-09-11}

\textbf{Author:} Matthew J. Colbrook, Department of Applied Mathematics
and Theoretical Physics, University of Cambridge. \textbf{Independent
Codex-agent review: PASS for the exact target.}

The ordered theorem transfers ordinary relative Frobenius residual error
with no loss for the larger monotone subhomogeneous function class. Put
\(B=\widehat A_k\). Since \(0\preceq B\preceq A\),
\(\|A-B\|_F^2=\|A\|_F^2-\|B\|_F^2-2\operatorname{tr}(B(A-B))\le\|A\|_F^2-\|B\|_F^2\).
Thus the original stronger trace-deficit premise implies the proved
residual premise, establishing the exact canonical conclusion. All
specified eigenbasis choices and zero-tail cases are covered.

The complete target is resolved. Its former difficulty rating is
historical; the original mathematical statement and source evidence are
retained below.

\textbf{Primary reference:}
\href{https://github.com/ajt60gaibb/OpenProblemsInNLA/blob/main/references/colbrook-transfer-2026-09-11/manuscripts/02_frobenius_function_transfer.pdf}{complete
authored PDF},
\href{https://github.com/ajt60gaibb/OpenProblemsInNLA/blob/main/references/colbrook-transfer-2026-09-11/manuscripts/02_frobenius_function_transfer.tex}{standalone
TeX}, \textbf{Theorem 1.1, with the trace-deficit reduction above}.
\href{https://github.com/ajt60gaibb/OpenProblemsInNLA/blob/main/references/colbrook-transfer-2026-09-11/verification/reviews/RA-09-review.md}{Independent
proof review} ·
\href{https://github.com/ajt60gaibb/OpenProblemsInNLA/blob/main/references/colbrook-transfer-2026-09-11/README.md}{Authorship
and submission record}. Verification is independent agent review, not
external human peer review or formal certification.

Let \(n\ge2\), \(1\le k<n\), and
\(A,\widehat A\in\mathbb R^{n\times n}\) be symmetric with
\(A\succeq\widehat A\succeq0\). Let \(f:[0,\infty)\to[0,\infty)\) be
continuous, concave, and nondecreasing. For a positive semidefinite
matrix \(X=\sum_{i=1}^n\lambda_iq_iq_i^T\) with decreasing eigenvalues
and orthonormal eigenvectors, write

\[
X_k=\sum_{i=1}^k\lambda_iq_iq_i^T,
\quad f(X)=\sum_{i=1}^nf(\lambda_i)q_iq_i^T,
\quad f(X)_k=\sum_{i=1}^kf(\lambda_i)q_iq_i^T.
\]

Use the same eigenvectors for the two truncations, with any choices
inside repeated eigenspaces. Is the implication

\[
\|A\|_F^2-\|\widehat A_k\|_F^2
\le(1+\varepsilon)\|A-A_k\|_F^2
\]

\[
\Longrightarrow\quad
\|f(A)-f(\widehat A)_k\|_F^2
\le(1+\varepsilon)\|f(A)-f(A)_k\|_F^2
\]

valid for every such input, every eigendecomposition choice, and every
\(\varepsilon\ge0\)? Here \(\|M\|_F^2=\sum_{i,j}|M_{ij}|^2\). The
premise is intentionally the difference of squared norms, not merely a
relative-error bound for \(\widehat A_k\).

The problem asks whether a reusable approximation to \(A\) provides
near-optimal Frobenius approximation to a larger class of matrix
functions.

\newpage
\subsection{References}\label{references}

\begin{enumerate}
\def\labelenumi{\arabic{enumi}.}
\tightlist
\item
  D. Persson, R. A. Meyer, and C. Musco,
  \href{https://arxiv.org/html/2311.14023v2}{Algorithm-agnostic low-rank
  approximation of operator monotone matrix functions}, \emph{SIAM
  Journal on Matrix Analysis and Applications} 46 (2025), 1--21. Table
  1, Frobenius/concave/ordered cell and its asterisk; Theorem 2.5
  supplies the stronger premise being extended, and §5 explicitly leaves
  the concave case open. Equation (2) specifies truncation.
\item
  D. Persson, T. Chen, and C. Musco,
  \href{https://arxiv.org/abs/2502.01888}{Randomized block-Krylov
  subspace methods for low-rank approximation of matrix functions},
  \emph{Linear Algebra and its Applications} 741 (2026), 32--65. This
  related work analyzes particular randomized Krylov constructions.
\end{enumerate}

\subsection{Status check --- 2026-09-08}\label{status-check-2026-09-08}

Checked the source's latest listed arXiv v2 (2024-07-04), 2025 journal
record, Table 1, Theorem 2.5, and §5. Searches combining the title,
``concave'', ``Frobenius'', ``counterexample'', ``funNyström'', and 2026
found no resolution. The related 2026 Krylov paper does not establish
the universal ordered-pair implication. No explicit reaffirmation later
than the 2025 publication was located. The source's nuclear-norm
counterexample and its unordered Frobenius example do not answer this
target. Unlike the spectral-error question, this statement concerns
squared Frobenius error and a stronger premise; these are distinct
explicitly marked open cells in the source table.

\subsection{Audit --- 2026-09-10}\label{audit-2026-09-10}

Rechecked \href{https://arxiv.org/html/2311.14023v2}{Table 1 and Theorem
2.5}: operator-monotone functions satisfy this stronger-premise
implication, while general concave functions remain open.
Frobenius/concavity and later matrix-function searches found no
resolution. The difference-of-squared-norms premise is essential.
\end{document}
